\documentclass[12pt,a4paper]{article}

% Packages
\usepackage[margin=1in]{geometry}
\usepackage{amsmath,amssymb,amsfonts}
\usepackage{graphicx}
\usepackage{booktabs}
\usepackage{multirow}
\usepackage{array}
\usepackage[table]{xcolor}
\usepackage{hyperref}
\usepackage{enumitem}
\usepackage{algorithm}
\usepackage{algpseudocode}
\usepackage{caption}
\usepackage{listings}
\usepackage{fancyhdr}
\usepackage{titlesec}
\usepackage{float}

% Colour Palette
\definecolor{darkblue}{RGB}{0,51,102}
\definecolor{lightblue}{RGB}{173,216,230}
\definecolor{darkgreen}{RGB}{0,100,0}
\definecolor{alertred}{RGB}{180,0,0}
\definecolor{codegray}{RGB}{245,245,245}
\definecolor{codeteal}{RGB}{0,128,128}

% Hyperref Setup
\hypersetup{
    colorlinks=true,
    linkcolor=darkblue,
    citecolor=darkgreen,
    urlcolor=codeteal,
    pdftitle={Aegis: Stochastic Ensemble Extension Report},
    pdfauthor={Jayant Kumar (230101050)},
}

% Header / Footer
\pagestyle{fancy}
\fancyhf{}
\rhead{\textcolor{darkblue}{\small Aegis Extension Report}}
\lhead{\textcolor{darkblue}{\small 230101050}}
\cfoot{\thepage}
\renewcommand{\headrulewidth}{0.4pt}

% Section Title Formatting
\titleformat{\section}{\large\bfseries\color{darkblue}}{\thesection}{1em}{}[\titlerule]
\titleformat{\subsection}{\normalsize\bfseries\color{darkblue}}{\thesubsection}{1em}{}

% Code Listing Style
\lstset{
    backgroundcolor=\color{codegray},
    basicstyle=\ttfamily\footnotesize,
    keywordstyle=\color{darkblue}\bfseries,
    commentstyle=\color{darkgreen},
    stringstyle=\color{alertred},
    breaklines=true,
    frame=single,
    rulecolor=\color{lightblue},
    captionpos=b,
    numbers=left,
    numberstyle=\tiny\color{gray},
    language=Python
}

% Custom Commands
\newcommand{\asr}{\textsc{ASR}}
\newcommand{\cda}{\textsc{CDA}}
\newcommand{\tbt}{\textsc{TBT}}
\newcommand{\sdn}{\textsc{SDN}}
\newcommand{\ic}{\textsc{IC}}

%=============================================================================
\begin{document}

% Title Page
\begin{titlepage}
    \centering
    \vspace*{1.5cm}
    {\Huge\bfseries\color{darkblue} Aegis: Mitigating Targeted\\[0.4em]
    Bit-Trojan Attacks on Neural Networks\par}
    \vspace{0.5cm}
    \rule{\linewidth}{2pt}
    \vspace{0.3cm}
    {\Large\itshape Original Framework Review and\\[0.2em]
    Stochastic Ensemble Extension\par}
    \vspace{1cm}
    \rule{\linewidth}{0.5pt}
    \vspace{1cm}
    {\large\bfseries Submitted By\par}
    \vspace{0.3cm}
    {\large Jayant Kumar\par}
    {\large Roll No.: 230101050\par}
    \vspace{1cm}
    {\large\bfseries Course\par}
    \vspace{0.3cm}
    {\large Computer Science --- Research Extension Project\par}
    \vspace{1cm}
    {\large\bfseries Date\par}
    \vspace{0.3cm}
    {\large August 2026\par}
    \vspace{2cm}
\end{titlepage}

\tableofcontents
\newpage

%=============================================================================
\section{Introduction}
\label{sec:intro}

As deep neural networks (DNNs) are increasingly deployed on resource-constrained edge hardware, their weights are stored in unprotected DRAM. Recent work has demonstrated that hardware-level \textit{Bit-Flip Attacks} (BFAs)~\cite{bfa2019} can precisely flip individual bits in DRAM using the Rowhammer exploit, allowing an adversary to either catastrophically degrade a model's performance or, more alarmingly, inject a covert backdoor that activates only when a specific trigger pattern is present in the input. The second class of attacks is formalized as \textit{Targeted Bit-Trojan} (\tbt{}) attacks~\cite{tbt2020}.

\textbf{Aegis}~\cite{aegis} is a recently proposed defense framework that leverages the \textit{Shallow-Deep Network} (\sdn{}) architecture~\cite{sdn} to mitigate such attacks. The core idea is that a single DNN can be augmented with multiple \textit{Internal Classifiers} (\ic{}s) at intermediate layers, and inference can be \emph{randomly routed} through a dynamic subset of these classifiers. Because an attacker who flips a single weight cannot predict which \ic{}s will be selected at inference time, the attack is structurally less effective.

This report has two main parts:
\begin{enumerate}
    \item \textbf{A review of the original Aegis framework}: its architecture, defense mechanism, and reported experimental results.
    \item \textbf{A description and evaluation of a novel extension}: a \textit{Self-Healing Stochastic Ensemble} that adds cryptographic weight-checksum verification on top of Aegis's randomized routing, enabling the system to detect compromised \ic{}s and automatically exclude them from the voting process.
\end{enumerate}

%=============================================================================
\section{Background}
\label{sec:background}

\subsection{Hardware Bit-Flip Attacks (BFA)}

Rowhammer-based Bit-Flip Attacks exploit electrical interference between adjacent DRAM rows to flip specific bits at will. In the context of DNNs, this translates to modifying the binary representation of a quantized model weight without requiring any software-level access. Since quantized weights are stored as compact integer types (e.g., 8-bit fixed-point), a single bit flip in the most significant bit can change a weight's value dramatically.

\subsection{Targeted Bit-Trojan (TBT) Attack}

The \tbt{} attack~\cite{tbt2020} is a precision BFA variant that implants a backdoor. The attack proceeds in two stages:

\begin{enumerate}
    \item \textbf{Neuron Gradient Ranking (NGR)}: The attacker performs a forward and backward pass on a representative data batch to identify the weight bits whose flipping will maximally push the model's output towards a target class for a specific trigger pattern.
    \item \textbf{Trojan Insertion}: Using FGSM-style gradient ascent, an optimal trigger image patch is crafted. Simultaneously, the identified bits in the victim model's weights are flipped to create a conditioned backdoor.
\end{enumerate}

After a successful \tbt{} attack, the model behaves normally on clean inputs but classifies any trigger-patched input as the target class, achieving a high \textit{Attack Success Rate} (\asr{}).

\subsection{Shallow-Deep Networks (SDN)}

\sdn{}s~\cite{sdn} are standard CNNs augmented with \ic{}s at the output of each residual block. Each \ic{} consists of a small sub-network (global average pooling + a quantized linear layer) that can produce a final classification. During vanilla \sdn{} inference, the network exits at the first \ic{} whose prediction confidence exceeds a threshold $\tau$ (typically $0.95$), saving computation for easy inputs.

%=============================================================================
\section{Aegis: The Original Framework}
\label{sec:aegis}

\subsection{Architecture}

Aegis is built on top of a quantized ResNet-32 \sdn{}. The key architectural parameters are:
\begin{itemize}
    \item \textbf{Backbone}: ResNet-32 with quantized weights (\texttt{quan\_Conv2d}, \texttt{quan\_Linear}).
    \item \textbf{Internal Classifiers}: 15 \ic{}s placed after each residual block + 1 final classifier = 16 total branches.
    \item \textbf{Datasets}: CIFAR-10, CIFAR-100, STL-10, Tiny-ImageNet.
    \item \textbf{Quantization}: 8-bit fixed-point for both weights and activations.
\end{itemize}

\subsection{Defense Mechanism: Randomized Dynamic Early-Exit}

Aegis's defense operates at inference time:
\begin{enumerate}
    \item Sort all \ic{}s by their clean validation accuracy.
    \item For each inference, randomly select a subset of $C$ high-performing branches.
    \item Evaluate them sequentially; exit at the first branch whose confidence $> \tau = 0.95$.
\end{enumerate}

The randomness of branch selection is the core security mechanism: an attacker who flips bits in one specific \ic{} cannot guarantee that \ic{} will be selected for any given inference.

\subsection{Training Procedure}

Aegis employs a two-stage training pipeline:
\begin{enumerate}
    \item \textbf{Base Training}: The backbone and all \ic{}s are jointly trained for 200 epochs using SGD with cosine annealing.
    \item \textbf{Branch Fine-Tuning}: Only the \ic{} linear layers are fine-tuned at a smaller learning rate to maximize per-branch accuracy.
\end{enumerate}

\subsection{Original Published Results}

Table~\ref{tab:original_aegis} reproduces key results from the original Aegis paper~\cite{aegis} for ResNet-32.

\begin{table}[H]
    \centering
    \caption{\asr{} (\%) of \tbt{} attacks on ResNet-32 from the original Aegis paper. Lower \asr{} is better.}
    \label{tab:original_aegis}
    \begin{tabular}{llccc}
        \toprule
        \textbf{Dataset} & \textbf{Attack} & \textbf{Base Model} & \textbf{SDN} & \textbf{Aegis} \\
        \midrule
        \multirow{2}{*}{CIFAR-10}   & Non-Adaptive & 72.3\% & 22.1\% & 13.0\% \\
                                     & Adaptive     & 70.7\% & 37.2\% & 31.1\% \\
        \midrule
        \multirow{2}{*}{CIFAR-100}  & Non-Adaptive & 89.4\% & 58.3\% & 31.9\% \\
                                     & Adaptive     & 95.8\% & 79.3\% & 49.7\% \\
        \midrule
        \multirow{2}{*}{STL-10}     & Non-Adaptive & ---    & ---    & 13.0\% \\
                                     & Adaptive     & 100.0\% & 35.0\% & 31.8\% \\
        \bottomrule
    \end{tabular}
\end{table}

%=============================================================================
\section{Extension: Self-Healing Stochastic Ensemble}
\label{sec:extension}

\subsection{Motivation}

The original Aegis framework has two key weaknesses:
\begin{enumerate}
    \item \textbf{No explicit detection}: Aegis \emph{probabilistically} reduces \asr{}, but never explicitly detects whether any \ic{} has been compromised. If the randomly selected \ic{} happens to be the trojanized one, the attack still succeeds.
    \item \textbf{Single-branch decision}: Only one branch makes the final decision (the first to exceed $\tau$). A single overconfident or trojanized branch can unilaterally dictate the outcome.
\end{enumerate}

This extension introduces two complementary mechanisms:
\begin{itemize}
    \item \textbf{SHA-256 Checksum Integrity Verification}: each \ic{}'s weights are hashed at deployment; any bit flip is detected with probability 1.
    \item \textbf{Confidence-Weighted Ensemble Voting}: all $C$ sampled (verified) branches vote simultaneously via a weighted average.
\end{itemize}

\subsection{Checksum-Based IC Integrity Verification}

\subsubsection{Rationale}

\tbt{} is specifically designed to be invisible to accuracy-based monitors: the model performs normally on clean data. A weight-level hash, however, changes the moment \emph{any} bit is flipped, regardless of the attack's behavioral effect. Table~\ref{tab:detector_comparison} summarizes the comparison.

\begin{table}[H]
    \centering
    \caption{Comparison of detection methods against BFA and TBT attacks.}
    \label{tab:detector_comparison}
    \begin{tabular}{lcc}
        \toprule
        \textbf{Detector Type} & \textbf{Catches BFA?} & \textbf{Catches TBT?} \\
        \midrule
        Accuracy/confidence monitoring & Yes & \textbf{No} (accuracy preserved by design) \\
        Weight checksum (full coverage) & Yes & \textbf{Yes} (any bit flip changes the hash) \\
        Weight checksum (partial only)  & Usually & Only if the flipped layer is covered \\
        \bottomrule
    \end{tabular}
\end{table}

\subsubsection{Algorithm}

\begin{algorithm}[H]
\caption{Checksum Verification and Self-Healing Vote Exclusion}
\label{alg:checksum}
\begin{algorithmic}[1]
\State \textbf{At Deployment:}
\For{each \ic{} $i \in \{0, \ldots, N-1\}$}
    \State $h_i^* \leftarrow \text{SHA256}(\text{state\_dict}(\ic{}_i))$ \Comment{golden checksum}
\EndFor
\State \textbf{At Inference:}
\State $\mathit{trusted} \leftarrow \{\}$
\For{each \ic{} $i$ in randomly sampled subset $S$}
    \State $h_i \leftarrow \text{SHA256}(\text{state\_dict}(\ic{}_i))$
    \If{$h_i = h_i^*$}
        \State $\mathit{trusted} \leftarrow \mathit{trusted} \cup \{i\}$ \Comment{IC intact}
    \Else
        \State \textbf{Alert}: \ic{} $i$ compromised; exclude permanently from pool
    \EndIf
\EndFor
\State Compute ensemble prediction over $\mathit{trusted}$ only (Eq.~\ref{eq:ensemble})
\end{algorithmic}
\end{algorithm}

\subsubsection{Properties}

\begin{itemize}
    \item \textbf{Detection probability}: 100\% for any single bit flip (collision resistance of SHA-256).
    \item \textbf{Computational cost}: Hashing a ResNet-32 \ic{}'s weights ($\sim$30\,KB) takes $<1$\,ms.
    \item \textbf{No retraining required}: Checksum computed once post-training.
\end{itemize}

\subsection{Stochastic Confidence-Weighted Voting}

Once integrity-verified branches are identified, the final prediction is:

\begin{equation}
    P_{\text{ensemble}} = \frac{\displaystyle\sum_{i \in \mathit{trusted}} W_i \cdot P_i}{\displaystyle\sum_{i \in \mathit{trusted}} W_i}
    \label{eq:ensemble}
\end{equation}

where $P_i = \text{Softmax}(\text{logits}_i)$ is the probability distribution from branch $i$, and $W_i = \max(P_i)$ is its peak confidence (the voting weight). The final class label is:

\begin{equation}
    \hat{y} = \arg\max\left(P_{\text{ensemble}}\right)
    \label{eq:prediction}
\end{equation}

Advantages over single-branch early-exit:
\begin{itemize}
    \item No single branch unilaterally determines the outcome.
    \item Higher-confidence branches receive proportionally more influence.
    \item The fixed $\tau = 0.95$ hyperparameter is eliminated.
\end{itemize}

\subsection{Self-Healing Behavior}

When a compromised \ic{} is detected:
\begin{enumerate}
    \item It is excluded from the current inference's ensemble vote.
    \item Its index is permanently removed from the valid sampling pool.
\end{enumerate}
This creates a self-healing system: as compromised branches are progressively excluded, clean accuracy \emph{improves} because only the highest-performing uncompromised branches remain active.

\subsection{Implementation}

The extension was implemented in Python/PyTorch as new scripts layered on top of the existing Aegis codebase:
\begin{itemize}
    \item \texttt{TBT\_nonadaptive\_ensemble.py} --- Non-adaptive \tbt{} with ensemble defense (CIFAR-10 \& STL-10)
    \item \texttt{TBT\_adaptive\_ensemble.py} --- Adaptive \tbt{} with ensemble defense (CIFAR-10 \& STL-10)
    \item SLURM submission scripts with \texttt{\$TMPDIR} staging for cluster execution
\end{itemize}

\begin{lstlisting}[caption={Core checksum generation and IC verification}, label={lst:checksum}]
import hashlib, torch

def compute_ic_checksum(model: torch.nn.Module) -> str:
    hasher = hashlib.sha256()
    for name, param in model.state_dict().items():
        hasher.update(param.detach().cpu().numpy().tobytes())
    return hasher.hexdigest()

# --- At deployment ---
golden = {i: compute_ic_checksum(ic) for i, ic in enumerate(ics)}

# --- At inference ---
trusted = []
for i, ic in enumerate(sampled_ics):
    if compute_ic_checksum(ic) == golden[i]:
        trusted.append(i)
        print(f"[OK]    IC {i} verified -- model intact.")
    else:
        print(f"[ALERT] IC {i} COMPROMISED -- excluded from vote.")
\end{lstlisting}

%=============================================================================
\section{Experimental Setup}
\label{sec:setup}

\subsection{Models and Datasets}

All experiments use a \textbf{ResNet-32} \sdn{} with 16 total branches. Two datasets are evaluated:
\begin{itemize}
    \item \textbf{CIFAR-10}: 50{,}000 train / 10{,}000 test, $32{\times}32$ pixels, 10 classes.
    \item \textbf{STL-10}: 5{,}000 train / 8{,}000 test, $96{\times}96$ pixels, 10 classes.
\end{itemize}

\subsection{Attacks Evaluated}

\begin{description}[style=nextline]
    \item[Non-Adaptive \tbt{}] Trigger and bit-flips optimized against the undefended base model. The attacker is unaware of stochastic routing or checksum.
    \item[Adaptive \tbt{}] White-box attacker with full knowledge of the \sdn{} architecture. The trigger is simultaneously optimized to corrupt \emph{all} branch classifiers' feature representations.
\end{description}

\subsection{Metrics}

\begin{itemize}
    \item \textbf{\cda{}} (Clean Data Accuracy): classification accuracy on unperturbed inputs.
    \item \textbf{\asr{}} (Attack Success Rate): fraction of triggered inputs classified as the target class.
\end{itemize}

\subsection{Hardware}

NVIDIA P100 (16\,GB) and A100 (40\,GB) GPUs on a SLURM-managed cluster.

%=============================================================================
\section{Results}
\label{sec:results}

\subsection{CIFAR-10 --- Non-Adaptive Attack}

\begin{table}[H]
    \centering
    \caption{Non-Adaptive \tbt{} on CIFAR-10 (ResNet-32). $\downarrow$\asr{} and $\uparrow$\cda{} are better.}
    \label{tab:cifar10_nonadaptive}
    \begin{tabular}{lccc}
        \toprule
        \textbf{Metric} & \textbf{Baseline Single Model} & \textbf{Aegis (Original)} & \textbf{Aegis + Ensemble (Ours)} \\
        \midrule
        \cda{} (Undefended) & 88.69\% & 64.47\% & 90.09\% \\
        \cda{} (Defended)   & N/A     & N/A      & \textbf{90.30\%} \\
        \midrule
        \asr{} (Undefended) & 17.82\% & 14.76\% & 15.31\% \\
        \asr{} (Defended)   & N/A     & 13.00\% & \textbf{12.62\%} \\
        \bottomrule
    \end{tabular}
\end{table}

Observations:
\begin{itemize}
    \item The ensemble alone boosts \cda{} by $+8.08\%$ over vanilla Aegis routing (64.47\% $\to$ 90.09\%).
    \item After checksum defense fires: \asr{} falls further from 15.31\% to 12.62\%, and \cda{} rises to 90.30\%.
\end{itemize}

\subsection{CIFAR-10 --- Adaptive Attack}

\begin{table}[H]
    \centering
    \caption{Adaptive \tbt{} on CIFAR-10 (ResNet-32). $\downarrow$\asr{} and $\uparrow$\cda{} are better.}
    \label{tab:cifar10_adaptive}
    \begin{tabular}{lccc}
        \toprule
        \textbf{Metric} & \textbf{Baseline Single Model} & \textbf{Aegis (Original)} & \textbf{Aegis + Ensemble (Ours)} \\
        \midrule
        \cda{} (Undefended) & 83.41\% & 67.31\% & 87.36\% \\
        \cda{} (Defended)   & N/A     & N/A      & \textbf{90.66\%} \\
        \midrule
        \asr{} (Undefended) & 38.50\% & 46.06\% & 40.41\% \\
        \asr{} (Defended)   & N/A     & 31.10\% & \textbf{9.79\%}  \\
        \bottomrule
    \end{tabular}
\end{table}

Observations:
\begin{itemize}
    \item \textbf{Critical result}: The adaptive attacker achieves 40.41\% \asr{} on the undefended ensemble. After checksum verification triggers and compromised \ic{}s are excluded, \asr{} drops to \textbf{9.79\%} --- equivalent to random guessing across 10 classes.
    \item \cda{} \emph{improves} from 87.36\% to \textbf{90.66\%} after self-healing because the accuracy-degrading trojanized sub-networks are removed.
    \item The baseline single model's clean accuracy degrades to 83.41\%; the ensemble recovers to 90.66\% (+7.25\%).
\end{itemize}

\subsection{CIFAR-10 --- Comprehensive Summary}

\begin{table}[H]
    \centering
    \caption{Full CIFAR-10 comparison across all model variants.}
    \label{tab:cifar10_full}
    \begin{tabular}{llcc}
        \toprule
        \textbf{Model} & \textbf{Attack} & \textbf{\cda{} (Post-Defense)} & \textbf{\asr{} (Post-Defense)} \\
        \midrule
        Baseline Single         & Non-Adaptive & 88.69\% & 17.82\% \\
        Aegis (Original)        & Non-Adaptive & 64.47\% & 13.00\% \\
        \rowcolor{lightblue}
        Aegis + Ensemble (Ours) & Non-Adaptive & \textbf{90.30\%} & \textbf{12.62\%} \\
        \midrule
        Baseline Single         & Adaptive & 83.41\% & 38.50\% \\
        Aegis (Original)        & Adaptive & 67.31\% & 31.10\% \\
        \rowcolor{lightblue}
        Aegis + Ensemble (Ours) & Adaptive & \textbf{90.66\%} & \textbf{9.79\%} \\
        \bottomrule
    \end{tabular}
\end{table}

\subsection{STL-10 Results}

\begin{table}[H]
    \centering
    \caption{TBT attack results on STL-10 (ResNet-32, $96\times96$ inputs, 5{,}000 training samples).}
    \label{tab:stl10}
    \begin{tabular}{llccc}
        \toprule
        \textbf{Attack} & \textbf{Metric} & \textbf{Baseline Aegis} & \textbf{Ensemble (Undefended)} & \textbf{Ensemble (Defended)} \\
        \midrule
        \multirow{2}{*}{Non-Adaptive}
          & \cda{} & ---     & 73.89\% & 74.11\% \\
          & \asr{} & 13.00\% & 12.21\% & \textbf{10.21\%} \\
        \midrule
        \multirow{2}{*}{Adaptive}
          & \cda{} & ---     & 73.70\% & 74.19\% \\
          & \asr{} & 31.00\% & 28.25\% & \textbf{13.91\%} \\
        \bottomrule
    \end{tabular}
\end{table}

Observations:
\begin{itemize}
    \item Adaptive \asr{}: from 31.00\% (baseline Aegis) to \textbf{13.91\%} with the ensemble defense.
    \item Non-adaptive \asr{}: from 13.00\% to \textbf{10.21\%}.
    \item Clean accuracy is marginally improved by self-healing in both attack scenarios.
    \item The relative improvement pattern is \textbf{consistent across both CIFAR-10 and STL-10}, confirming generalizability of the approach to larger-resolution, smaller-dataset settings.
\end{itemize}

\subsection{Why Does \asr{} Initially Rise Under Adaptive Attack?}

In Table~\ref{tab:cifar10_adaptive}, the undefended ensemble shows a \emph{higher} \asr{} (40.41\%) than original Aegis (31.10\%). This is expected: the adaptive attacker explicitly knows the ensemble structure and crafts the trigger to simultaneously corrupt all branch classifiers, making the undefended ensemble more vulnerable. However, this aggressive multi-branch modification also means the trojanized ICs produce clearly distinguishable weight modifications that the checksum catches with certainty. Table~\ref{tab:trajectory} tracks this through the full pipeline.

\begin{table}[H]
    \centering
    \caption{\asr{} trajectory through defense stages under the adaptive attack (CIFAR-10).}
    \label{tab:trajectory}
    \begin{tabular}{lc}
        \toprule
        \textbf{Defense Stage} & \textbf{\asr{}} \\
        \midrule
        No defense (baseline single model) & 38.50\% \\
        Aegis randomized routing (original) & 31.10\% \\
        Our ensemble, no checksum applied & 40.41\% \\
        \rowcolor{lightblue}
        \textbf{Our ensemble + checksum (self-healed)} & \textbf{9.79\%} \\
        \bottomrule
    \end{tabular}
\end{table}

%=============================================================================
\section{Discussion}
\label{sec:discussion}

\subsection{Comparison with Original Aegis}

\begin{enumerate}
    \item \textbf{Deterministic vs.\ probabilistic}: Aegis \emph{reduces the probability} of hitting a compromised \ic{}. Our extension \emph{detects} the compromise and eliminates it.
    \item \textbf{Accuracy recovery}: The original Aegis accepts accuracy degradation caused by trojanized weights. Self-healing restores accuracy to above-attack levels by routing only through verified clean branches.
\end{enumerate}

\subsection{Computational Overhead}

\begin{itemize}
    \item \textbf{Checksum}: SHA-256 over a ResNet-32 \ic{} (${\sim}30$\,KB) takes $<1$\,ms. For 16 ICs total, $<16$\,ms --- negligible vs.\ forward-pass time.
    \item \textbf{Ensemble forward passes}: $C=5$ branches per inference vs.\ $\sim1$ in Aegis early-exit. A $5\times$ overhead on the sampled subset, but still far cheaper than full 16-branch evaluation.
\end{itemize}

\subsection{Limitations}

\begin{enumerate}
    \item \textbf{Secure checksum storage}: Golden hashes must reside in tamper-resistant storage (TEE, secure flash). If the attacker can corrupt both the model and the stored checksums, detection fails.
    \item \textbf{STL-10 model quality}: The STL-10 model achieves only 73.85\% accuracy due to the small (5{,}000 sample) training set. The original Aegis paper likely used semi-supervised pre-training on STL-10's 100{,}000 unlabeled images, explaining the gap to their reported numbers.
    \item \textbf{Hash collision}: A sufficiently determined attacker with knowledge of the checksum scheme could theoretically craft a weight modification that produces an identical SHA-256 hash. This is computationally infeasible with SHA-256 but a theoretical concern for weaker checksums.
\end{enumerate}

%=============================================================================
\section{Conclusion}
\label{sec:conclusion}

This work presented and evaluated a \textit{Self-Healing Stochastic Ensemble} extension to the Aegis defense framework, combining:
\begin{itemize}
    \item \textbf{SHA-256 weight checksums} for deterministic, behavior-independent detection of any bit modification.
    \item \textbf{Confidence-weighted stochastic voting} across a randomly sampled subset of verified \ic{}s.
    \item \textbf{Automatic self-healing} that permanently excludes compromised \ic{}s and recovers clean accuracy.
\end{itemize}

\noindent\textbf{Key results on CIFAR-10 (ResNet-32):}
\begin{itemize}
    \item Adaptive \asr{}: $31.10\%$ (Aegis) $\rightarrow$ \textbf{9.79\%} (near-random guessing across 10 classes).
    \item Non-adaptive \asr{}: $13.00\%$ (Aegis) $\rightarrow$ \textbf{12.62\%}.
    \item Clean accuracy under adaptive attack: $67.31\%$ (Aegis) $\rightarrow$ \textbf{90.66\%}.
\end{itemize}

\noindent\textbf{Key results on STL-10 (ResNet-32):}
\begin{itemize}
    \item Adaptive \asr{}: $31.00\%$ (Aegis) $\rightarrow$ \textbf{13.91\%}.
    \item Non-adaptive \asr{}: $13.00\%$ (Aegis) $\rightarrow$ \textbf{10.21\%}.
\end{itemize}

The self-healing ensemble is a training-free, inference-time defense that generalizes across datasets and attack types, representing a principled and practical advancement over the base Aegis framework.

%=============================================================================
\section*{Acknowledgements}
The author thanks the original Aegis authors for making their codebase publicly available, and the cluster administrators for GPU access.

%=============================================================================
\begin{thebibliography}{9}

\bibitem{bfa2019}
A.~S.~Rakin, Z.~He, and D.~Fan,
``Bit-Flip Attack: Crushing Neural Network with Progressive Bit Search,''
\textit{ICCV}, 2019.

\bibitem{tbt2020}
A.~S.~Rakin, Z.~He, J.~Li, F.~Yao, C.~Chakrabarti, and D.~Fan,
``TBT: Targeted Neural Network Attack with Bit Trojan,''
\textit{CVPR}, 2020.

\bibitem{aegis}
J.~Li, A.~S.~Rakin, X.~Xiong, L.~Chang, Z.~He, and D.~Fan,
``Defending Bit-Flip Attack through DNN Weight Reconstruction,''
\textit{DAC}, 2021.

\bibitem{sdn}
B.~Kaya, T.~M.~Dumlu, and A.~Yazici,
``Shallow-Deep Networks: Understanding and Mitigating Network Overthinking,''
\textit{ICML}, 2019.

\bibitem{rowhammer}
Y.~Kim et~al.,
``Flipping Bits in Memory Without Accessing Them: An Experimental Study of DRAM Disturbance Errors,''
\textit{ISCA}, 2014.

\bibitem{strip}
Y.~Gao, C.~Xu, D.~Wang, S.~Chen, D.~C.~Ranasinghe, and S.~Nepal,
``STRIP: A Defence Against Trojan Attacks on Deep Neural Networks,''
\textit{ACSAC}, 2019.

\end{thebibliography}

\end{document}
